\section{Analysis}
\label{sec:analysis}

\paragraph{Why two models lack the gap.}
The no-gap models differ. \mbox{GPT-5 submits} $92\%$ of expected entries and has \bsexact{} essentially equal to \bsrecon{} ($0.508$ vs.\ $0.500$). \mbox{Grok-4.3 submits} only $40\%$ of expected entries, so its balance sheet is internally aligned but incomplete (Appendix Table~\ref{tab:no-gap-mechanism}).

\paragraph{Where the gap concentrates.}
On the anchor model, \bsexact{} drops from $0.88$ at L1 to $0.00$ at L5 while \bsrecon{} declines more gently; the \aggGap{} peaks near L2--L3 at $\sim$50\,pp (Figure~\ref{fig:difficulty-trend}). Ledger feedback narrows the gap at every level. Industry and concept slices in Appendix~\ref{appx:per-concept} show that retail and SaaS remain entry-limited.

\paragraph{Error taxonomy.}
On $120$ standard records, document-linking failures dominate the anchor model: $92.5\%$ of records and $981$ entries, compared with $235$ omissions, $139$ extras, $37$ arithmetic errors, and only $14$ account-selection errors. The hard part is grounding plausible entries to the right evidence. For detailed model-failure slices, see Appendix Figure~\ref{fig:failure-slices}.
